\documentclass[12pt]{ca-capstone}
\geometry{margin=25mm}
\usepackage{setspace}
\setstretch{1.1}
\capstonerefs{crandall-pomerance,hps,msu-ent}
\begin{document}
\capstonetitle{Тесты простоты и псевдопростые числа}
\begin{center}\textbf{Теоретическая справка для преподавателя. Не выдавать студентам.}\end{center}

\section{Ферма и Миллер--Рабин}
Для простого $p$ и $\gcd(a,p)=1$ выполняется $a^{p-1}=1\pmod p$. Нарушение даёт свидетеля
непростоты, выполнение не доказывает простоту: основание может быть лжецом. Число Кармайкла
составное, но проходит тест для всех взаимно простых с ним оснований.

В Миллере--Рабине $n-1=d2^s$, $d$ нечётно. Для простого $n$ последовательность
$a^{d},a^{2d},\ldots,a^{2^{s}d}$ либо начинается с $1$, либо содержит $-1$ перед первым $1$.
Отсутствие этих форм даёт сильного свидетеля. Сильные псевдопростые к выбранному основанию
нужны для проверки ложного принятия.

\section{Соловей--Штрассен и Якоби}
Для нечётного простого $n$ выполняется $a^{(n-1)/2}\equiv\legendre{a}{n}\pmod n$. Символ Якоби
обобщает символ Лежандра и вычисляется без факторизации: выделение двоек использует класс
$n\bmod8$, перестановка --квадратичный закон взаимности, после чего $a$ заменяется на
$n\bmod a$. Несовпадение значения с $\pm1$ или нулевой Якоби при единичном НОД выдаёт составное
число. Эйлеровы псевдопростые проходят это сравнение при выбранном основании.

\section{Корсельт и число лжецов}
Составное $n$ является числом Кармайкла тогда и только тогда, когда оно свободно от квадратов и
\[
 p-1\mid n-1\quad\text{для каждого простого }p\mid n.
\]
Все делимости могут выполняться и при $p^2\mid n$, но такое $n$ не Кармайклово; эти примеры
обязательны. Для свободного от квадратов $n$ число взаимно простых с $n$ оснований-лжецов равно
\[
 L=\prod_{p\mid n}\gcd(n-1,p-1),\qquad \varphi(n)=\prod_{p\mid n}(p-1).
\]
Формула следует из КТО и числа решений $x^{n-1}=1$ в $\GF{p}^*$.

\section{Люк--Лемер и приведение Мерсенна}
Для $M_p=2^p-1$ при нечётном простом $p$ полагают $s_0=4$,
$s_{i+1}=s_i^2-2\bmod M_p$. Тест утверждает, что $M_p$ просто тогда и только тогда, когда
$s_{p-2}=0$; $p=2$ рассматривается отдельно.

Приведение использует $H2^p+L\equiv H+L\pmod{M_p}$. Повторяют его до получения значения не
больше $M_p$; значение ровно $M_p$ заменяется нулём. Обычное деление на $M_p$ скрывает существо
требования и запрещено. Тест включает как простые, так и составные $M_p$ (например, $M_{11}$).

\clearpage
\section{Вердикты и контрольный список}
Wilson даёт точный критерий $(n-1)!\equiv-1\pmod n$, но факториальная стоимость делает его
только малым оракулом. Вероятностные тесты возвращают \texttt{probable\_prime} или
\texttt{composite}. Основания должны принадлежать допустимому диапазону $1<a<n-1$; малые простые
и кратные $n$ основания не должны давать ложный отрицательный результат.

Для каждого свидетеля независимо повторяются НОД, степень, последовательность Миллера--Рабина
или символ Якоби. Проверяются фиксированные списки Кармайкла, сильных и эйлеровых псевдопростых,
не-свободные от квадратов примеры Корсельта, повторяемость оснований и обнуление финальной
невязки Люка--Лемера в точности для простых чисел Мерсенна.

\references
\end{document}
